\section{Seguimiento de ejemplos}

\subsection{Ejemplo 1}

Tomemos el arreglo de monedas.

\[
M = [20, 5, 10, 85, 30]
\]

El algoritmo \textit{greedy} del TP1 nos hubiera dado el siguente resultado, tomando la mayor entre las monedas disponibles.

\begin{itemize}
    \item \textbf{Monedas de Sophia:} $[30, 20, 10] \implies \text{Ganancia} = 30 + 20 + 10 = 60$
    \item \textbf{Monedas de Mateo:} $[85, 5] \implies \text{Ganancia} = 85 + 5 = 90$
\end{itemize}

Si comparamos las ganancias obtenidas ($60$ vs. $90$), se observa claramente que el enfoque \textit{greedy} no es optimo, ya que la ganancia de Mateo resulta ser mayor a la de Sophia.

Por otro lado, con nuestro algoritmo, realizamos los siguentes pasos.

\begin{enumerate}
    \item Creamos la matriz optima de $N \times N$ 
    \begin{table}[h!]
    \centering
    \begin{tabular}{cccccc}
    \toprule
    $i \setminus j$ & 1 & 2 & 3 & 4 & 5 \\
    \midrule
    1 & -- & -- & -- & -- & -- \\
    2 & -- & -- & -- & -- & -- \\
    3 & -- & -- & -- & -- & -- \\
    4 & -- & -- & -- & -- & -- \\
    5 & -- & -- & -- & -- & -- \\
    \bottomrule
    \end{tabular}
    \end{table}

    \item Llenamos los casos base
    \begin{table}[h!]
    \centering
    \begin{tabular}{cccccc}
    \toprule
    $i \setminus j$ & 1 & 2 & 3 & 4 & 5 \\
    \midrule
    1 & 20 & 20 & -- & -- & -- \\
    2 & -- & 5  & 10 & -- & -- \\
    3 & -- & -- & 10 & 85 & -- \\
    4 & -- & -- & -- & 85 & 85 \\
    5 & -- & -- & -- & -- & 30 \\
    \bottomrule
    \end{tabular}
    \end{table}

    \item Terminamos de completar la matriz, segun nos indica el algoritmo
    \begin{table}[h!]
    \centering
    \begin{tabular}{cccccc}
    \toprule
    $i \setminus j$ & 1 & 2 & 3 & 4 & 5 \\
    \midrule
    1 & 20 & 20 & 25 & 95 & 110 \\
    2 & -- & 5  & 10 & 90 & 90  \\
    3 & -- & -- & 10 & 85 & 40  \\
    4 & -- & -- & -- & 85 & 85  \\
    5 & -- & -- & -- & -- & 30  \\
    \bottomrule
    \end{tabular}
    \end{table}
\end{enumerate}

Vemos entonces que nuestra nueva ganancia maxima, encontrada en la esquina superior derecha de nuestra tabla, es de $110$, mayor a la de nuestro algoritmo \textit{greedy}.Finalmente para obtener la secuencia de monedas que deberia elegir Sophia a partir de la matriz $OPT$ ya calculada para $M = [20, 5, 10, 85, 30]$, el recorrido recursivo desde el estado inicial $(i=1, j=5)$ es el siguente.

\begin{enumerate}
    \item \textbf{Estado inicial} $(i=1, j=5) \implies \text{Rango } [20, 5, 10, 85, 30]$:
    \begin{itemize}
        \item Si Sophia elige $M[1] = 20$: Mateo ve $M[2]=5$ y $M[5]=30$. Elige $30$. Queda el subproblema $OPT[2][4] = 90$.
        \item Si Sophia elige $M[5] = 30$: Mateo ve $M[1]=20$ y $M[4]=85$. Elige $85$. Queda el subproblema $OPT[1][3] = 25$.
        \item \textbf{Decisión}: Sophia elige $20$.
        \item \textbf{Moneda guardada}: $M[1] = 20$.
        \item \textbf{Siguiente estado}: $(i=2, j=4)$.
    \end{itemize}

    \item \textbf{Estado} $(i=2, j=4) \implies \text{Rango } [5, 10, 85]$:
    \begin{itemize}
        \item Si Sophia elige $M[2] = 5$: Mateo ve $M[3]=10$ y $M[4]=85$. Elige $85$. Queda $OPT[3][3] = 10$.
        \item Si Sophia elige $M[4] = 85$: Mateo ve $M[2]=5$ y $M[3]=10$. Elige $10$. Queda $OPT[2][2] = 5$.
        \item \textbf{Decisión}: Sophia elige $85$.
        \item \textbf{Moneda guardada}: $M[4] = 85$.
        \item \textbf{Siguiente estado}: $(i=2, j=3)$.
    \end{itemize}

    \item \textbf{Estado} $(i=2, j=3) \implies \text{Rango } [5, 10]$ (Caso base de 2 elementos):
    \begin{itemize}
        \item Sophia elige la moneda máxima entre $M[2]=5$ y $M[3]=10$.
        \item \textbf{Decisión}: Sophia elige $10$.
        \item \textbf{Moneda guardada}: $M[3] = 10$.
    \end{itemize}
\end{enumerate}

\textbf{Resultado final del seguimiento:}
\begin{itemize}
    \item \textbf{Monedas elegidas por Sophia ($s$):} $[20, 85, 10]$.
    \item \textbf{Monedas elegidas por Mateo:} $[30, 5]$.
\end{itemize}